\SourceLeaf{293}%
so ist die folgende Umformung für eine direkte, strenge Lösung am günstigsten. Die beiden ersten Formeln bleiben dabei dieselben wie in der ersten Methode.

Vierte Methode
\[
\begin{aligned}
\operatorname{tang}s&=\cos T\operatorname{tang}r\\
\operatorname{tang}\lambda&=\frac{\operatorname{tang}T\sin s}{\cos(S-s)}\\
\operatorname{tang}t&=\sin T\sin r\operatorname{tang}(S-s)\\
\sin\tau&=\sin T\operatorname{tang}\tfrac{1}{2}r\sin s\\
\sin\sigma&=\operatorname{tang}t\operatorname{tang}\tfrac{1}{2}\lambda\cos(S-s)\\
S'&=S-s-\sigma\\
T'&=T-t-\tau
\end{aligned}
\]
\indent Diese vierte Methode lässt an Genauigkeit nichts zu wünschen übrig. Bei der unmittelbaren Rechnung in dieser Form muss man jedoch etwas mühsam interpolieren, um die kleinen Bögen aus den Logarithmen der Tangenten oder Sinus zu bestimmen. Diesen Nachteil kann man leicht umgehen, indem man die trigonometrischen Funktionen in Reihen entwickelt. Damit kann man die Rechnungen ohne Genauigkeitsverlust mit den Logarithmen der Zahlen durchführen. Es genügt, hier nur die wichtigsten Schritte dieser Umformung anzugeben.

Es sei
\[
\begin{aligned}
r\cos T&=s^0\\
r\sin T&=v
\end{aligned}
\]
Bezeichnet man dann zur Abkürzung die Größe eines Bogens von einer Sekunde in Teilen des Radius oder den Bruch \(\tfrac{\pi}{648000}\) mit \(\rho\) und betrachtet \(r\) als Größe erster Ordnung, so gilt bis zu Größen fünfter Ordnung ausschließlich
\[
s=s^0(1+\tfrac{1}{3}\rho\rho rr-\tfrac{1}{3}\rho\rho s^0s^0)=s^0(1+\tfrac{1}{3}\rho\rho vv)
\]
Setzt man dann außerdem
\[
\begin{aligned}
v\operatorname{tang}(S-s)&=t^0\\
\frac{v}{\cos(S-s)}&=\lambda^0
\end{aligned}
\]
so erhält man
